\documentclass{zhvt-classic}

% 使用本地源流明体（GenRyuMin2 TW）按文件路径加载，不依赖系统字体名
% 版框（内框／界栏／书口侧栏／鱼尾）改用与版心印章同一印泥棕色：
%   印章图片 yshs.png 实色 = RGB(135,68,52) = HTML 874434
\usepackage{xcolor}
\usepackage{tikz}
% sealbrown 双注册：xcolor 供 TikZ/pgf，L3 供 cls 版框
\definecolor{sealbrown}{HTML}{874434}
\ExplSyntaxOn
\color_set:nnn { sealbrown } { HTML } { 874434 }
\ExplSyntaxOff

\zhvtset{
  grid color       = sealbrown,
  main font        = GenRyuMin2TW-SB,
  main font path   = C:/Users/Administrator/texmf/fonts/opentype/public/GenRyuMin2TW/,
jiazhu font      = GenRyuMin2TW-M,
  jiazhu font path = C:/Users/Administrator/texmf/fonts/opentype/public/GenRyuMin2TW/,
  % 页码框容纳 4 字：正常页「九十七」+前置空 1 格=4 行恰好；
  % 100+ 页「一百零一」共 5 行超出，末字「一/二/三」被框裁掉（2026-10-04 实测），故放宽到 5
  shukou page number chars = 5,
  % 书口书名框高从默认 6 字减到 4 字，把省出的 2 字高拨给 chapter 框，
  % 使 6 字书名+卷名放进 9 字高的 chapter 框后不再压线（2026-10-05）
  shukou title chars = 4,
}

\title{史记}[鲁仲连邹阳列传]
\author{汉·司马迁}
\maker{闲叟漫记·重製}


% ===== 索隐述赞样式（2026-11-08 四次定稿 · 愚大人亲手改定参数）=====
% ⚠ 竖排前提：自定义字体家族必须带 RawFeature={vertical:+vert}；
%   simhei.ttf 加 +vhal 会整体加载失败静默回退，只写 +vert；simfang 同理。
\newCJKfontfamily{\shuzanfont}[Path=C:/Windows/Fonts/,RawFeature={vertical:+vert}]{simhei.ttf}
\newCJKfontfamily{\shuzanbody}[Path=C:/Windows/Fonts/,RawFeature={vertical:+vert}]{simfang.ttf}
% 题头：152x30pt 圆角 6pt 棕底白字；baseline (0,2) 使整框落入界行格子不压版框线；
% 白字节点 x 不补偿、y +8.2pt（像素级墨迹实测 dx/dy≈0，L/R 留白约 20pt）。
\newcommand{\shuzanhead}[1]{%
  \begin{tikzpicture}[baseline={(0,2)}]
    \fill[sealbrown, rounded corners=6pt]
      (-76pt,-15pt) rectangle (76pt,15pt);
    \node[text=white, inner sep=0pt, anchor=center] at (0pt,8.2pt)
      {\shuzanfont\fontsize{23}{23}\selectfont #1};
  \end{tikzpicture}%
}
% ★ 修正书口卷名位置：chapter 框起始位置改回 cls 默认值，避免全局影响。
%   只在每个卷首页通过 \juan 命令里的空命令微调该页位置（安全，不会跨页污染）。
\ExplSyntaxOn
\AtBeginDocument{
  \dim_set:Nn \l__zhvt_shukou_chapter_title_start_dim
    { \l__zhvt_shukou_title_height_dim + \l__zhvt_font_size_dim * 2 }
  \dim_set:Nn \l__zhvt_shukou_page_number_start_dim
    { \l__zhvt_shukou_chapter_title_start_dim + \l__zhvt_shukou_chapter_title_height_dim }
}
\ExplSyntaxOff

\begin{document}

% ★★ 版心（书口）排布覆写：把书名放到鱼尾之下（2026-10-04）
%
%   cls 默认版心自上而下：书名（title 静态框，在鱼尾之上）→ 鱼尾 → 篇名（chapter 动态框）→ 页码
%   本模板改为：          鱼尾（位置不动）→ 书名 → 篇名 → 页码分割线 → 页码
%
%   手法：把 cls 放在 title 静态框里的书名，改放进 chapter 动态框（鱼尾之下）；
%         title 静态框留空。书名取 \l__zhvt_title_str，即 \title{...} 的第一参数，
%         所以换书时只改 \title，不必在此重复写书名。
%
%   ⚠ 三条必须遵守的写法（2026-10-04 实测）：
%     ① 汉字必须用 \hbox{...} 包住。flowfram 会把框内容装进一个宽 14.4mm（= 一个 line_sep
%        = 书口栏宽）的纵向 minipage 并按该宽度折行，而整页靠 PDF /Rotate 90 显示、
%        框内文字本身不变向 —— 不包 \hbox 时「围炉夜话」4 字会被逐字折成 4 行，
%        在版面上就成了横排 4 字并横向溢出版框。
%     ② 竖排下**一个 \hbox 就是一整列**。书名与篇名要同处一列，必须放进同一个 \hbox；
%        分成两个 \hbox 时，后一个会被当成新的一列、甩到隔壁栏去
%        （实测「序」跑到书口栏外侧 14.4mm 的正文列上）。中间的 \quad 在框内是横向空白，
%        显示出来正好是书名与篇名之间的纵向间隔。
%     ③ ★ 内容盒要写成 **与框等宽的定宽盒** `\hbox to \l__zhvt_shukou_chapter_title_height_dim{…\hfil}`。
%        实测：chapter 框是把内容盒沿纵向**居中**摆放的，内容宽度一变，书名就跟着漂：
%        自然宽度的 \hbox（书名 38.4mm）落在 x=157.2mm，加篇名后（57.6mm）落在 x=152.4mm，
%        目录页（无篇名）与正文页的书名就差半格。做成等宽盒后，无论有没有篇名，
%        内容固定从框顶起排，实测两页均落在 x=142.8mm，不再漂。
\ExplSyntaxOn
% 把书口 chapter 框（书名+篇名）起始位置下移 1 字高，避免 4 字卷名顶到鱼尾
% ⚠ cls 的 \__zhvt_calculate_dimensions: 在 \AtBeginDocument 里执行，会重算该 dim；
%   必须把覆写挂到 begin document 之后，否则会被二次覆盖（2026-10-05 实测）。
% (旧方案已移到导言区，直接修改 chapter 框起始位置)
% 卷标题命令 #1=书口篇名(鱼尾) #2=正文卷名（含卷序，卷名内空格用 ~）
%
% ★★ 目录页两项修订（2026-10-05，愚大人指出）★★
%   ① 页码起点统一对齐：
%      原状 —— 目录项写成 \par\hspace{4em}卷名\hspace{4em}页码，页码起点 = 卷名实宽 + 4em。
%      卷名实宽随字数变：「卷一~梁惠王上」6 字 +~、「卷七~离娄上」5 字 +~，
%      差 1 个字位 → 卷七~卷十 4 条的页码整体比其余 10 条高 27.21pt（一个字），看着上下不齐。
%      改法 —— 卷名补全角空格（宽 = 1 字 = 27.215pt）补足到 **7 字符位**再写进目录，
%      14 条页码起点遂全部落在同一条竖线上（起点对齐；末字随页码字数自然参差）。
%      ⚠ 只对目录里的副本 \l_tmpc_tl 补，正文卷名仍用 \l_tmpa_tl 原样 —— 正文像素不受影响。
%   ② 卷名整体上移半字位：前导 \hspace 由 4em 改 3.5em。
%      实测原状目录块在版框内偏下（天头留白 139.6pt / 地脚留白 110.7pt），
%      上移 0.5em = 13.6pt 后上下基本均分，视觉居中。
\NewDocumentCommand { \juan } { mm } {
	\clearpage
	\tl_set:Nn \l__zhvt_chapter_title_tl { #1 }
	\tl_set:Nn \l_tmpa_tl { #2 }
	\tl_set_eq:NN \l_tmpc_tl \l_tmpa_tl
	\int_step_inline:nn { 7 - \str_count:n { #2 } }
	  { \tl_put_right:Nn \l_tmpc_tl { 　 } }
	\tl_set:Nx \l_tmpb_tl { \zhnum { page } }
	\tl_put_left:Nn \l_tmpb_tl { \hspace { 4em } }
	\tl_put_left:Nn \l_tmpb_tl { \l_tmpc_tl }
	\tl_put_left:Nn \l_tmpb_tl { \par \hspace { 3.5em } }
	\addtocontents { toc } { \l_tmpb_tl }
	\skip_horizontal:n { \l__zhvt_font_size_dim * 4 }
	\l_tmpa_tl
	\par
	\thispagestyle { empty }
}
\cs_set_protected:Npn \__zhvt_set_shukou_content: {
	% ★ 鱼尾上方留白（title 静态框）嵌入闲叟钤印：
	%   整页 /Rotate 90，故此处把图用 \rotatebox{90} 转正为竖章；
	%   尺寸卡在版心条幅内（书内约 7mm 宽 × 33.8mm 高，留余量不压鱼尾）。
	\setstaticcontents * { title } {
		\vbox_to_ht:nn { \l__zhvt_baseline_skip_dim } {
			\vskip 14.2pt
			\hbox_to_wd:nn { \l__zhvt_shukou_title_height_dim } {
				\hfil
				\rotatebox [ origin = c ] { 90 } { \includegraphics [ width = 7mm , height = 33.8mm ] { yshs.png } }
				\hfil
			}
		}
	}
% ★ 版心（书口）统一改法（2026-11-16 愚大人定）：18pt＋框内居中＋书口文字同版框色
%   ① 原 27.2pt 时 7 字长篇名（如「廉颇蔺相如列传」）末字溢出 chapter 框、压住页码；
%      改 \fontsize{18pt}{18pt}，前后 \hfil 令字串在框内居中（不随篇名字数漂移）。
%   ② \vbox{\vskip 8.2pt} 定量下移，抵消缩字号后栏宽方向 8.2pt 上偏（像素实测 y 中心 543.3 = 原版）。
%      ⚠ \raisebox / \vspace* 在 flowfram 框内无效；换字号须同步换补偿：16pt→10.1、19pt→7.4。
%   ③ chapter 与 page~number 均置 \color{sealbrown}（HTML 874434），与版框／界行同色；正文仍黑。
%   ⚠ 页码字位、钤印、鱼尾、界行、正文一律未动；五部合订本同此口径。
	\setdynamiccontents * { chapter }
	{ \color{sealbrown} \vbox { \vskip 8.2pt \hbox to \l__zhvt_shukou_chapter_title_height_dim
		{ \hfil \fontsize{18pt}{18pt}\selectfont \l__zhvt_title_str \quad \l__zhvt_chapter_title_tl \hfil } } }
	\setdynamiccontents * { page~number } { \color{sealbrown} \quad \zhnum { page } }
}
\__zhvt_set_shukou_content:
\ExplSyntaxOff

\maketitle

% 单篇，无多卷目录，不渲染目录页

\juan{鲁仲连邹阳列传}{卷八十三~鲁仲连邹阳列传第二十三}
\hspace*{2em}鲁仲连者。齐人也。好奇伟俶傥之画策。而不肯仕宦任职。好持高节。游於赵。\\
\hspace*{2em}赵孝成王时。而秦王使白起破赵长平之军前後四十馀万。秦兵遂东围邯郸。赵王恐。诸侯之救兵莫敢击秦军。魏安釐王使将军晋鄙救赵。畏秦。止於荡阴不进。魏王使客将军新垣衍间入邯郸。因平原君谓赵王曰。﹁秦所为急围赵者。前与齐湣王争彊为帝。已而复归帝。今齐已益弱。方今唯秦雄天下。此非必贪邯郸。其意欲复求为帝。赵诚发使尊秦昭王为帝。秦必喜。罢兵去。﹂平原君犹预未有所决。\\
\hspace*{2em}此时鲁仲连適游赵。会秦围赵。闻魏将欲令赵尊秦为帝。乃见平原君曰。﹁事将柰何。﹂平原君曰。﹁胜也何敢言事。前亡四十万之众於外。今又内围邯郸而不能去。魏王使客将军新垣衍令赵帝秦。今其人在是。胜也何敢言事。﹂鲁仲连曰。﹁吾始以君为天下之贤公子也。吾乃今然後知君非天下之贤公子也。梁客新垣衍安在。吾请为君责而归之。﹂平原君曰。﹁胜请为绍介而见之於先生。﹂平原君遂见新垣衍曰。﹁东国有鲁仲连先生者。今其人在此。胜请为绍介。交之於将军。﹂新垣衍曰。﹁吾闻鲁仲连先生。齐国之高士也。衍人臣也。使事有职。吾不原见鲁仲连先生。﹂平原君曰。﹁胜既已泄之矣。﹂新垣衍许诺。\\
\hspace*{2em}鲁连见新垣衍而无言。新垣衍曰。﹁吾视居此围城之中者。皆有求於平原君者也。今吾观先生之玉貌。非有求於平原君者也。曷为久居此围城之中而不去。﹂鲁仲连曰。﹁世以鲍焦为无从颂而死者。皆非也。众人不知。则为一身。彼秦者。弃礼义而上首功之国也。权使其士。虏使其民。彼即肆然而为帝。过而为政於天下。则连有蹈东海而死耳。吾不忍为之民也。所为见将军者。欲以助赵也。﹂\\
\hspace*{2em}新垣衍曰。﹁先生助之将柰何。﹂鲁连曰。﹁吾将使梁及燕助之。齐。楚则固助之矣。﹂新垣衍曰。﹁燕则吾请以从矣。若乃梁者。则吾乃梁人也。先生恶能使梁助之。﹂鲁连曰。﹁梁未睹秦称帝之害故耳。使梁睹秦称帝之害。则必助赵矣。﹂\\
\hspace*{2em}新垣衍曰。﹁秦称帝之害何如。﹂鲁连曰。﹁昔者齐威王尝为仁义矣。率天下诸侯而朝周。周贫且微。诸侯莫朝。而齐独朝之。居岁馀。周烈王崩。齐後往。周怒。赴於齐曰。﹃天崩地坼。天子下席。东籓之臣因齐後至。则斮。﹄齐威王勃然怒曰。﹃叱嗟。而母婢也。﹄卒为天下笑。故生则朝周。死则叱之。诚不忍其求也。彼天子固然。其无足怪。﹂\\
\hspace*{2em}新垣衍曰。﹁先生独不见夫仆乎。十人而从一人者。宁力不胜而智不若邪。畏之也。﹂鲁仲连曰。﹁呜呼。梁之比於秦若仆邪。﹂新垣衍曰。﹁然。﹂鲁仲连曰。﹁吾将使秦王烹醢梁王。﹂新垣衍怏然不悦。曰。﹁噫嘻。亦太甚矣先生之言也。先生又恶能使秦王烹醢梁王。﹂鲁仲鲁曰。﹁固也。吾将言之。昔者九侯。鄂侯。文王。纣之三公也。九侯有子而好。献之於纣。纣以为恶。醢九侯。鄂侯争之彊。辩之疾。故脯鄂侯。文王闻之。喟然而叹。故拘之牖里之库百日。欲令之死。曷为与人俱称王。卒就脯醢之地。齐湣王之鲁。夷维子为执策而从。谓鲁人曰。﹃子将何以待吾君。﹄鲁人曰。﹃吾将以十太牢待子之君。﹄夷维子曰。﹃子安取礼而来吾君。彼吾君者。天子也。天子巡狩。诸侯辟舍。纳筦籥。摄衽抱机。视膳於堂下。天子已食。乃退而听朝也。﹄鲁人投其籥。不果纳。不得入於鲁。将之薛。假途於邹。当是时。邹君死。湣王欲入吊。夷维子谓邹之孤曰。﹃天子吊。主人必将倍殡棺。设北面於南方。然后天子南面吊也。﹄邹之群臣曰。﹃必若此。吾将伏剑而死。﹄固不敢入於邹。邹。鲁之臣。生则不得事养。死则不得赙襚。然且欲行天子之礼於邹。鲁。邹。鲁之臣不果纳。今秦万乘之国也。梁亦万乘之国也。俱据万乘之国。各有称王之名。睹其一战而胜。欲从而帝之。是使三晋之大臣不如邹。鲁之仆妾也。且秦无已而帝。则且变易诸侯之大臣。彼将夺其所不肖而与其所贤。夺其所憎而与其所爱。彼又将使其子女谗妾为诸侯妃姬。处梁之宫。梁王安得晏然而已乎。而将军又何以得故宠乎。﹂\\
\hspace*{2em}於是新垣衍起。再拜谢曰。﹁始以先生为庸人。吾乃今日知先生为天下之士也。吾请出。不敢复言帝秦。﹂秦将闻之。为卻军五十里。適会魏公子无忌夺晋鄙军以救赵。击秦军。秦军遂引而去。\\
\hspace*{2em}於是平原君欲封鲁连。鲁连辞让者三。终不肯受。平原君乃置酒。酒酣起前。以千金为鲁连寿。鲁连笑曰。﹁所贵於天下之士者。为人排患释难解纷乱而无取也。即有取者。是商贾之事也。而连不忍为也。﹂遂辞平原君而去。终身不复见。\\
\hspace*{2em}其後二十馀年。燕将攻下聊城。聊城人或谗之燕。燕将惧诛。因保守聊城。不敢归。齐田单攻聊城岁馀。士卒多死而聊城不下。鲁连乃为书。约之矢以射城中。遗燕将。书曰。\\
\hspace*{2em}吾闻之。智者不倍时而弃利。勇士不却死而灭名。忠臣不先身而後君。今公行一朝之忿。不顾燕王之无臣。非忠也。杀身亡聊城。而威不信於齐。非勇也。功败名灭。後世无称焉。非智也。三者世主不臣。说士不载。故智者不再计。勇士不怯死。今死生荣辱。贵贱尊卑。此时不再至。原公详计而无与俗同。\\
\hspace*{2em}且楚攻齐之南阳。魏攻平陆。而齐无南面之心。以为亡南阳之害小。不如得济北之利大。故定计审处之。今秦人下兵。魏不敢东面。衡秦之势成。楚国之形危。齐弃南阳。断右壤。定济北。计犹且为之也。且夫齐之必决於聊城。公勿再计。今楚魏交退於齐。而燕救不至。以全齐之兵。无天下之规。与聊城共据期年之敝。则臣见公之不能得也。且燕国大乱。君臣失计。上下迷惑。栗腹以十万之众五折於外。以万乘之国被围於赵。壤削主困。为天下僇笑。国敝而祸多。民无所归心。今公又以敝聊之民距全齐之兵。是墨翟之守也。食人炊骨。士无反外之心。是孙膑之兵也。能见於天下。虽然。为公计者。不如全车甲以报於燕。车甲全而归燕。燕王必喜。身全而归於国。士民如见父母。交游攘臂而议於世。功业可明。上辅孤主以制群臣。下养百姓以资说士。矫国更俗。功名可立也。亡意亦捐燕弃世。东游於齐乎。裂地定封。富比乎陶。卫。世世称孤。与齐久存。又一计也。此两计者。显名厚实也。原公详计而审处一焉。\\
\hspace*{2em}且吾闻之。规小节者不能成荣名。恶小耻者不能立大功。昔者管夷吾射桓公中其钩。篡也。遗公子纠不能死。怯也。束缚桎梏。辱也。若此三行者。世主不臣而乡里不通。乡使管子幽囚而不出。身死而不反於齐。则亦名不免为辱人贱行矣。臧获且羞与之同名矣。况世俗乎。故管子不耻身在縲绁之中而耻天下之不治。不耻不死公子纠而耻威之不信於诸侯。故兼三行之过而为五霸首。名高天下而光烛邻国。曹子为鲁将。三战三北。而亡地五百里。乡使曹子计不反顾。议不还踵。刎颈而死。则亦名不免为败军禽将矣。曹子弃三北之耻。而退与鲁君计。桓公朝天下。会诸侯。曹子以一剑之任。枝桓公之心於坛坫之上。颜色不变。辞气不悖。三战之所亡一朝而复之。天下震动。诸侯惊骇。威加吴。越。若此二士者。非不能成小廉而行小节也。以为杀身亡躯。绝世灭後。功名不立。非智也。故去感忿之怨。立终身之名。弃忿悁之节。定累世之功。是以业与三王争流。而名与天壤相弊也。原公择一而行之。\\
\hspace*{2em}燕将见鲁连书。泣三日。犹豫不能自决。欲归燕。已有隙。恐诛。欲降齐。所杀虏於齐甚众。恐已降而後见辱。喟然叹曰。﹁与人刃我。宁自刃。﹂乃自杀。聊城乱。田单遂屠聊城。归而言鲁连。欲爵之。鲁连逃隐於海上。曰。﹁吾与富贵而诎於人。宁贫贱而轻世肆志焉。﹂\\
\hspace*{2em}邹阳者。齐人也。游於梁。与故吴人庄忌夫子。淮阴枚生之徒交。上书而介於羊胜。公孙诡之间。胜等嫉邹阳。恶之梁孝王。孝王怒。下之吏。将欲杀之。邹阳客游。以谗见禽。恐死而负累。乃从狱中上书曰。臣闻忠无不报。信不见疑。臣常以为然。徒虚语耳。昔者荆轲慕燕丹之义。白虹贯日。太子畏之。卫先生为秦画长平之事。太白蚀昴。而昭王疑之。夫精变天地而信不喻两主。岂不哀哉。今臣尽忠竭诚。毕议原知。左右不明。卒从吏讯。为世所疑。是使荆轲。卫先生复起。而燕。秦不悟也。原大王孰察之。\\
\hspace*{2em}昔卞和献宝。楚王刖之。李斯竭忠。胡亥极刑。是以箕子详狂。接舆辟世。恐遭此患也。原大王孰察卞和。李斯之意。而後楚王。胡亥之听。无使臣为箕子。接舆所笑。臣闻比干剖心。子胥鸱夷。臣始不信。乃今知之。原大王孰察。少加怜焉。\\
\hspace*{2em}谚曰。﹁有白头如新。倾盖如故。﹂何则。知与不知也。故昔樊於期逃秦之燕。藉荆轲首以奉丹之事。王奢去齐之魏。临城自刭以卻齐而存魏。夫王奢。樊於期非新於齐。秦而故於燕。魏也。所以去二国死两君者。行合於志而慕义无穷也。是以苏秦不信於天下。而为燕尾生。白圭战亡六城。为魏取中山。何则。诚有以相知也。苏秦相燕。燕人恶之於王。王按剑而怒。食以駃夬騠。白圭显於中山。中山人恶之魏文侯。文侯投之以夜光之璧。何则。两主二臣。剖心坼肝相信。岂移於浮辞哉。\\
\hspace*{2em}故女无美恶。入宫见妒。士无贤不肖。入朝见嫉。昔者司马喜髌脚於宋。卒相中山。范睢摺胁折齿於魏。卒为应侯。此二人者。皆信必然之画。捐朋党之私。挟孤独之位。故不能自免於嫉妒之人也。是以申徒狄自沈於河。徐衍负石入海。不容於世。义不苟取。比周於朝。以移主上之心。故百里奚乞食於路。缪公委之以政。甯戚饭牛车下。而桓公任之以国。此二人者。岂借宦於朝。假誉於左右。然後二主用之哉。感於心。合於行。亲於胶漆。昆弟不能离。岂惑於众口哉。故偏听生奸。独任成乱。昔者鲁听季孙之说而逐孔子。宋信子罕之计而囚墨翟。夫以孔。墨之辩。不能自免於谗谀。而二国以危。何则。众口铄金。积毁销骨也。是以秦用戎人由余而霸中国。齐用越人蒙而彊威。宣。此二国。岂拘於俗。牵於世。系阿偏之辞哉。公听并观。垂名当世。故意合则胡越为昆弟。由余。越人蒙是矣。不合。则骨肉出逐不收。硃。象。管。蔡是矣。今人主诚能用齐。秦之义。後宋。鲁之听。则五伯不足称。三王易为也。\\
\hspace*{2em}是以圣王觉寤。捐子之之心。而能不说於田常之贤。封比干之後。修孕妇之墓。故功业复就於天下。何则。欲善无厌也。夫晋文公亲其雠。彊霸诸侯。齐桓公用其仇。而一匡天下。何则。慈仁殷勤。诚加於心。不可以虚辞借也。\\
\hspace*{2em}至夫秦用商鞅之法。东弱韩。魏。兵彊天下。而卒车裂之。越用大夫种之谋。禽劲吴。霸中国。而卒诛其身。是以孙叔敖三去相而不悔。於陵子仲辞三公为人灌园。今人主诚能去骄泬之心。怀可报之意。披心腹。见情素。堕肝胆。施德厚。终与之穷达。无爱於士。则桀之狗可使吠尧。而蹠之客可使刺由。况因万乘之权。假圣王之资乎。然则荆轲之湛七族。要离之烧妻子。岂足道哉。\\
\hspace*{2em}臣闻明月之珠。夜光之璧。以闇投人於道路。人无不按剑相眄者。何则。无因而至前也。蟠木根柢。轮囷离诡。而为万乘器者。何则。以左右先为之容也。故无因至前。虽出随侯之珠。夜光之璧。犹结怨而不见德。故有人先谈。则以枯木朽株树功而不忘。今夫天下布衣穷居之士。身在贫贱。虽蒙尧。舜之术。挟伊。管之辩。怀龙逢。比干之意。欲尽忠当世之君。而素无根柢之容。虽竭精思。欲开忠信。辅人主之治。则人主必有按剑相眄之迹。是使布衣不得为枯木朽株之资也。\\
\hspace*{2em}是以圣王制世御俗。独化於陶钧之上。而不牵於卑乱之语。不夺於众多之口。故秦皇帝任中庶子蒙嘉之言。以信荆轲之说。而匕首窃发。周文王猎泾。渭。载吕尚而归。以王天下。故秦信左右而杀。周用乌集而王。何则。以其能越挛拘之语。驰域外之议。独观於昭旷之道也。\\
\hspace*{2em}今人主沈於谄谀之辞。牵於帷裳之制。使不羁之士与牛骥同皁。此鲍焦所以忿於世而不留富贵之乐也。\\
\hspace*{2em}臣闻盛饰入朝者不以利汙义。砥厉名号者不以欲伤行。故县名胜母而曾子不入。邑号朝歌而墨子回车。今欲使天下寥廓之士。摄於威重之权。主於位势之贵。故回面汙行以事谄谀之人而求亲近於左右。则士伏死堀穴岩之中耳。安肯有尽忠信而趋阙下者哉。\\
\hspace*{2em}书奏梁孝王。孝王使人出之。卒为上客。\\
\hspace*{2em}太史公曰。鲁连其指意虽不合大义。然余多其在布衣之位。荡然肆志。不诎於诸侯。谈说於当世。折卿相之权。邹阳辞虽不逊。然其比物连类。有足悲者。亦可谓抗直不桡矣。吾是以附之列传焉。\\

% ===== 索隐述赞（司马贞《史记索隐》）· 统一样式 =====
% 题头「索隐述赞曰」五字并入反白框；正文每四句一列，定宽盒 \hfil 居列中；
% \fontsize{24}{41} 列距对齐全书界行（121+41k）；\clearpage 使述赞独占一页。
\clearpage
\shuzanhead{索隐述赞曰}\linebreak
{\shuzanbody\fontsize{24}{41}\selectfont
\hbox to \hsize{\hfil 鲁连达士。高才远致。释难解纷。辞禄肆志。\hfil}\linebreak
\hbox to \hsize{\hfil 齐将挫辩。燕军沮气。邹子遇谗。见诋狱吏。\hfil}\linebreak
\hbox to \hsize{\hfil 慷慨献说。时王所器。\hfil}%
\par}
\end{document}
